% !TEX root = ../Projektdokumentation.tex
\section{Implementierung} 
\label{sec:Implementierung}
Dieser Abschnitt beschreibt, wie bei der Erstellung der Anwendung vorgegangen wurde. 
Zuerst wurden die Datenbanksturkturen und Anbindung implementiert. Danach wurde der Microservice erstellt.
Anschließend wurden die Message Broker Abhängigkeiten und der HTTP-Client für die Verbindung zum \acs{CRM} eingebaut.
Schließlich wurden die Datenbankanfragen bzw. Synchronisationshistorie und die Geschäftslogik entwickelt.   

\subsection{Implementierung der Datenstrukturen}
\label{sec:ImplementierungDatenstrukturen}
Die PostgreSQL Datenbank \textit{nightcrawler} wurde händisch über \textit{pgAdmin 4} angelegt und nutzt die \textit{libpq-15.dll} Bibliothek. 
Im Anschluss wurden mithilfe des Paketes \textit{\acs{EF-Core}}, innerhalb der Applikation Migrationen aus Beziehungen zwischen Entitäten erstellt und
auf die Datenbank angewandt. Ebenfalls wurden dadurch die Tabellen und Eigenschaften angelegt.


\subsection{Implementierung RabbitMQ}
\label{sec:ImplementierungRabbitMQ}
Für die Verbindung zwischen dem Microservice \textit{Nightcrawler} und dem Message Broker \textit{RabbitMQ} wurde mithilfe von 
Serviceregistrierung in der \acs{DI}-Klasse eine \textit{IConnectionFactory} als Singleton mit den Applikationseinstellungen registriert.
Durch die Registrierung als Singleton können alle Klassen, die eine Verbindung zu \textit{RabbitMQ} benötigen, mit Dependency Injection
dieselbe Instanz erzeugen. Dadurch müssen nicht alle Klassen die Einstellungen kennen, um die richtige \textit{RabbitMQ} Instanz anzusprechen.
\par\smallskip
Um die verschiedenen Queues zu initialisieren wurde ein QueueInitializer erstellt, welcher als \textit{IHostedService} implementiert wird,
um beim Start der Applikation die Queues zu erstellen bzw. zu versichern, dass sie vorhanden sind. Innerhalb dieses Services wird die Topologie
für \textit{RabbitMQ} verwaltet. Hier werden die \textit{MainQueues}, \textit{RetryQueues} und \textit{\ac{DLQ}'s} erstellt, die später für
den Eventverkehr und das Errorhandling benötigt werden.
\par\smallskip
Als letztes werden EventPublisher und EventConsumer benötigt, um Events senden und empfangen zu können.
Hierzu wurde im \textit{user-service} auf Infrastruktur-Ebene ein EventPublisher implementiert, welcher indirekt nach einem \acs{API}-Aufruf,
die zu synchronisierenden Daten als Event sendet. Der \textit{Nightcrawler} beinhaltet auf Presentations-Ebene den EventConsumer und empfängt die gesendeten Daten,
um diese dann zu verarbeiten und mit \textit{HubSpot} zu synchronisieren.
\par\smallskip
Screenshots der Anwendung in Bezug auf \textit{RabbitMQ} befinden sich im \Anhang{app:RabbitMQConnection}.


\subsection{Implementierung des HTTP-Clients}
\label{sec:ImplementierungHTTPClient}
Um nach dem Empfangen eines Events die Synchronisierung in Richtung \textit{HubSpot} anzustoßen wird ein HTTP-Client benötigt.
Dafür wurde der \textit{HubSpotApiClient} implementiert, welcher die zu sendenden Daten als Domänen-Model Parameter erhält, 
diese als Body setzt und damit einen HTTP request ausführt. Dieser Client ermöglicht POST und PATCH Anfragen bzw. das Erstellen und Bearbeiten von Kontaktdatensätzen in \textit{HubSpot}.
\par\smallskip
Damit die Anfragen die korrekte \textit{HubSpot} Instanz erreicht benötigt es eine \textit{BaseUrl}, also den ersten Teil einer \acs{URL}, 
wodurch die Instanz angesprochen wird. Desweiteren wird der Kontaktpfad benötigt, also der Ort, an dem die Kontaktdaten verändert werden können.
\textit{HubSpot} benötigt für ihre Schnittstellen eine Authorisierung, weshalb der Authorization Header mit Bearer und einem Zugriffstoken für diese HTTP-Anfragen
zusätzlich benötigt wird.
\par\smallskip
Um schließlich alle Anforderungen zu erfüllen wurde ein \textit{HubSpotRateLimiter} implementiert, der dafür sorgt, dass die von \textit{HubSpot}
festgelegten Anfragebegrenzungen nicht überschritten werden und somit keine Anfragen unnötig fehlschlagen oder verloren gehen.
\par\smallskip
Ausschnitte zum \textit{HubSpotApiClient} und \textit{HubSpotRateLimiter} sind im \Anhang{app:HubSpot} zu finden.


\subsection{Implementierung der Synchronisationshistorie}
\label{sec:ImplementierungSync}
Um die Synchronisationshistorie persistieren zu können wurden drei Abstraktionsebenen entwickelt, um auf allen Ebenen testbar zu bleiben
und die Aufteilung von Infrastruktur und Applikationsebene zu erhalten. Konkret handelt es sich dabei um \textit{DbOperations}, \textit{Repositories}
und Services jeweils pro, den bereits genannten \textit{SyncHistory} und \textit{UserContactState}.
\par\smallskip
Die \textit{DbOperations} befinden sich auf der niedrigsten Schicht und sie können einfache Abfragen auf Datenbankebene machen. 
Außerdem sind neben Testklassen, die einzigen Klassen, die den \textit{DbContext} kennen.
\par\smallskip
Bei den \textit{Repositories} handelt es sich, um Klassen, die \textit{DbOperations} implementieren und neben diesen Aufrufen Domänen-Models
zu Infrastruktur-Models konvertieren. Zudem befindet sich das Logging, Retry- und Errorhandling auf dieser Ebene. Für die Retrylogik wurde das NuGet Paket \textit{Polly}
verwendet, welches ermöglicht vorkonfigurierte \textit{ResiliencePipelines} zu implementieren, mit denen man fehlerhafte Aufrufe für eine bestimmte Anzahl
an Wiederholungen erneut ausführen kann.
\par\smallskip
Die Services beinhalten \textit{Repository}-Aufrufe, sowie minimale Logik, mit der \textit{Upserts} ermöglicht werden.
So überprüft man für die \textit{SyncHistory} zum Beispiel anhand der \textit{EventId}, ob es bereits eine Historie gibt. 
Wenn es keine gibt wird eine erstellt, andernfalls wird sie nur aktualisiert.
\par\smallskip
Siehe \Anhang{app:SyncHistory} für Code-Ausschnitte.


\subsection{Implementierung der Geschäftslogik}
\label{sec:ImplementierungGeschaeftslogik}
Für die Umsetzung der eigentlichen Geschäftslogik wurden drei weitere Services auf Applikationsebene implementiert, damit auch hier die Testbarkeit
und das \textit{single responsibility principle} eingehalten wird. Zwei Services für die Ressourcen \textit{PersonalInformation} und \textit{CommunicationPreference}
werden, je nach erhaltener Nachricht im EventConsumer instanziiert und starten den Synchronisationsprozess, der jeweilig deserialisierten Nachricht.
Desweiteren beinhalten diese Services die Logik der Erstellung von Synchroinsationshistorien.
\par\smallskip
Als letztes wurde der \textit{HubSpotContactService} entwickelt. Dieser beinhaltet den \textit{HubSpotApiClient} und führt die beiden genannten Ressourcen
ein \textit{mapping} durch, womit diese in eine für \textit{HubSpot} lesbare Form gebracht werden. Fakt ist, dass ein Kontakt nur erstellt werden soll, wenn
die Ressource \textit{PersonalInformation} das erste Mal synchronisiert wird. Aus diesem Grund wurden in dieser Klasse Methoden für das Erstellen und Aktualisieren
von Kontakten implementiert. Wenn die \textit{contactId}, also die Identifikationsnummer für einen \textit{HubSpot} Kontakt null beträgt, wird ein Kontakt erstellt.
Andernfalls wird der Kontakt mit der vorhandenen \textit{contactId} aktualisiert. Für den Fall, dass ein Kontakt mit der gegebenen E-Mail bereits in \textit{HubSpot}
existiert, wurde ein \textit{\acs{RegEx}} eingebaut, welcher anhand des Inhalts, im \textit{\acs{Json}} Format einer API Antwort mit Statuscode 409, die existierende
Indentifikationsnummer extrahiert und zurückgibt. Die \textit{contactId} wird schließlich in dem \textit{UserContactState} gespeichert, wodurch bei erneuter Datenänderung
des Nutzers, nurnoch die Daten aktualisiert werden.
\par\smallskip
Siehe \Anhang{app:Logic}
